\scanpage{19}
% L01–07
le foncteur $C\longrightarrow C'$ défini ainsi : donné $(G,\varphi)\in\mathrm{Ob}\,C$, on définit $\alpha$ comme l'invariant de $\varphi(\underline u)=u$, et $R$ comme le \deleted{fermé des} \unsure{sous foncteur|des foncteurs|faisceau} de $X$ (\inserted{plan projectif} enveloppe de la droite proj. $X/S$ définie par $G$) \deleted{fermé des} $H_1\setminus\{c\cup d\}$, qui est un torseur sous $U_{2-\alpha}\underset{\text{iso can}}{\sim}\deleted{\unsure{\mathrm{Aut}(G,\varphi)}}=\uline{\mathrm{Aut}}(X,\varphi)$

% L08–12
\ldots\ \unsure{Quant à|tenant à|tenant} : $\uline{\mathrm{Aut}}(\deleted{\unsure{Q}}X,\varphi)$, c'est le ss groupe \deleted{\uline{Corollaire} Soit en la \unsure{situation|relation}} de $U_{\beta\unsure{:=|=}2-\alpha}$ formé des $[\lambda]$ tels que
\[
\lambda(2+\lambda\beta)=0\quad\text{i.e.}\quad\deleted{\unsure{2\ne}}\ \lambda\mathbin{\ast_\beta}\lambda=1 .
\]
Donc

% L13–20, crossed by long diagonals
\deleted{1) Cas \unsure{$2\beta\ne0$|$2\rho\ne0$} \inserted{\unsure{$\beta=2-\alpha\ne0$|inversible}} \deleted{i.e. $\alpha\ne2$}, \deleted{i.e.} \inserted{i.e. $U_\beta\simeq$ un tore $\alpha\ne2)$ et $\beta\ne2$. \inserted{(car \unsure{inv.})} Alors} \deleted{\unsure{$U_\beta$ est tore. Alors}}}

\deleted{$\uline{\mathrm{Aut}}(X,\varphi)\simeq\deleted{\unsure{\mathbb Z_2}\ \unsure{\Leftrightarrow}\mu_2}\subset G_m$, \inserted{engendré par $\sigma$.} D'ailleurs l'image de $\deleted{\unsure{G_\alpha}}U_\alpha$ dans \deleted{Aut} $G\deleted{\unsure{(X)}}$ est l'unique tore maximal qui contient \deleted{\unsure{l'él.|l'image}} $u$, donc si $u$ est d'ordre fini pair $2n$, on a $\sigma=u^{\unsure{n|m}}$.}
% slanted
\margin{\deleted{\unsure{Dès que} \inserted{dans ces deux cas,} \unsure{dès qu'on pose|dès qu'on a} \unsure{on a} une représentation fidèle de $\mathbb D_{\unsure{\alpha|i|\infty}}\times\mathbb Z_2$ ($i$ pair) dont les \unsure{extensions|restrictions} à $\mathbb D_i^{\unsure{\alpha}}$ sont déter\unsure{.}}}

% L21–26
\deleted{2) Cas $\beta=0$, \deleted{$p\ne2$} \deleted{$2\ne0$} (i.e. $U_\beta\simeq G_a$ i.e. $\alpha=2$) \inserted{et \deleted{$p\ne2$} $2$ inv.} Alors, l'image de $U_\beta$ dans $G$ est l'unique \deleted{\unsure{un sous}} ssgroupe \deleted{unipotent} $\simeq G_a$ qui contient $u$. \deleted{\unsure{lorsque $p\ne$}} \deleted{\unsure{Si $2$ inv., toujours}} $\uline{\mathrm{Aut}}(X,\varphi)=\deleted{\unsure{\mu}}\ 1$}
% archive sheet number: 18
